\chapter{Conclusion}

This project investigates whether plain feedforward artificial neural networks can approximate the solution of the Lorenz-1960 ordinary differential equation system when trained only on high-accuracy numerical reference data. The motivation comes from a practical and academic tension. Classical numerical solvers such as Runge-Kutta methods provide dependable reference solutions, but they solve the system step by step on a fixed grid. Neural approaches can instead learn a continuous input-output mapping from time to the state variables, but recent work on differential equations often relies on physics-informed losses or operator-learning architectures \cite{matthews2024pinnde, babni2025error, lau2020oden}. This study therefore keeps a narrow scope: a data-driven feedforward ANN is trained to approximate $(x(t), y(t), z(t))$ for the Lorenz-1960 initial value problem, without adding ODE residual terms to the loss function.

The research design centers on a validated numerical baseline and a controlled architecture comparison. The Lorenz-1960 formulation, parameters, initial conditions, and time interval follow the system described in Chapter~2 from the PinnDE reference problem \cite{matthews2024pinnde}. The reference trajectory is generated and checked through two numerical solvers, a custom RK4 implementation and SciPy DOP853, before it is used as training and validation data. The ANN study is then divided into two phases. Phase~1 compares depth, width, and activation function while holding the optimizer fixed. Phase~2 compares Adam, L-BFGS, and SGD with momentum on the best Phase~1 architectures. This separation is intended to make the architecture and optimizer effects easier to interpret.

The expected contribution of the work is an empirical benchmark for plain ANN approximation of the Lorenz-1960 system. Existing studies show that neural networks can solve or approximate differential equations under several formulations \cite{mall2013comparison, okereke2021novel, dubey2024solving}, and physics-informed approaches have been applied to Lorenz-type problems \cite{matthews2024pinnde, aslam2024neuro}. The specific academic value of this project is its focus on a simpler question: how far a standard feedforward ANN can go when architecture choices are varied systematically and the training target is a verified numerical solution. The answer can help clarify whether better plain-ANN design is sufficient for this benchmark or whether physics-informed constraints remain necessary for higher accuracy and reliability.

At the current report stage, the supported contribution is the validated baseline design, the dataset and training protocol, and the planned evaluation framework. The report does not claim a final best ANN architecture unless the corresponding experiment records, metric tables, and plots are present in Chapter~4. This limitation is deliberate. A conclusion based on unsupported accuracy claims would weaken the study. Final claims about the best depth, width, activation function, or optimizer should be added only after the complete experimental evidence has been generated and reviewed.

After the planned comparison is completed, future work should extend the study in three directions. First, the best configuration should be tested under repeated random seeds to measure stability. Second, the trained model should be evaluated on denser time grids or additional initial conditions to assess generalization beyond the current trajectory. Third, the plain ANN results should be compared more directly with physics-informed and operator-learning baselines from the literature. These extensions would help determine whether the architecture selected in this project is only effective for one controlled trajectory or whether it can support broader neural approximation of nonlinear ODE systems.
